%% Abstract (configurado para língua inglesa)
\begin{resumo}[Abstract]
\begin{otherlanguage*}{english}
The accelerated growth of the Creator Economy and the consolidation of distance 
learning in Brazil have created a demand for educational content management tools 
that generic platforms such as Trello, Notion and Asana are unable to meet, as they 
lack native semantic structure for the educational domain. This paper describes the 
development of Creator Workbench, a web-based Content Operations (ContentOps) 
platform designed specifically to centralize and optimize the production pipeline for 
digital educational content creators. The methodology adopted is Scrum, with 
biweekly Sprints, a five-member team with distinct technical roles and a versioned 
repository on GitHub. The system architecture follows the decoupled Client-Server 
pattern, with a React.js front-end, a Java Spring Boot back-end, a MySQL relational 
database normalized to Third Normal Form and infrastructure provisioned on Amazon 
Web Services. The platform offers modules for idea management, a Kanban workflow 
with production states, a script editor, digital asset upload integrated with AWS S3, 
Python-based automations and curricular structuring into Courses, Modules and 
Lessons. Security is implemented through JWT authentication, BCrypt encryption and 
HTTPS communication, in compliance with the Brazilian General Data Protection Law 
(LGPD). The financial viability analysis demonstrates the sustainability of the SaaS 
model with three subscription plans, presenting an operating margin above 90\% from 
the seventh month in the realistic scenario. The results indicate that a tool with native 
semantics for the educational domain reduces the creator's operational fragmentation 
and creates the technical conditions for applying andragogical principles in the 
production of digital courses.

\vspace{\onelineskip}
\noindent
\textbf{Keywords}: Creator Economy. Content Operations. Educational Content Management. Distance Learning. SaaS.
\end{otherlanguage*}
\end{resumo}

%% Exemplo de resumo em francês
%\begin{resumo}[Résumé]
% \begin{otherlanguage*}{french}
%    Il s'agit d'un résumé en français.
% 
%   \textbf{Mots-clés}: latex. abntex. publication de textes.
% \end{otherlanguage*}
%\end{resumo}

%% Exemplo de resumo em Espanhol
%\begin{resumo}[Resumen]
% \begin{otherlanguage*}{spanish}
%   Este es el resumen en español.
%  
%   \textbf{Palabras clave}: latex. abntex. publicación de textos.
% \end{otherlanguage*}
%\end{resumo}
% ---